\[ %%%%%%%%%%%%%%%%%%%%%%%%%%%%%%%%%%%%%%%%%%%%%%%%%%%%%%%%%%%%%%%%%%%%%%%%%%%%%%%% \subsection{Benchmark Dynamical Systems} \label{sec:benchmark_systems} %%%%%%%%%%%%%%%%%%%%%%%%%%%%%%%%%%%%%%%%%%%%%%%%%%%%%%%%%%%%%%%%%%%%%%%%%%%%%%%% To evaluate the generality of PHYSGEN, experiments are conducted on a diverse set of nonlinear dynamical systems representing different levels of physical complexity. The selected benchmark suite includes: \begin{enumerate} \item low-dimensional chaotic dynamical systems; \item spatially distributed nonlinear partial differential equations; \item complex fluid dynamical systems. \end{enumerate} These benchmarks provide complementary evaluation scenarios for prediction accuracy, physical consistency, stability, and scalability. %%%%%%%%%%%%%%%%%%%%%%%%%%%%%%%%%%%%%%%%%%%%%%%%%%%%%%%%%%%%%%%%%%%%%%%%%%%%%%%% \subsubsection{Benchmark Selection Criteria} %%%%%%%%%%%%%%%%%%%%%%%%%%%%%%%%%%%%%%%%%%%%%%%%%%%%%%%%%%%%%%%%%%%%%%%%%%%%%%%% The benchmark systems are selected according to the following criteria: \begin{itemize} \item \textbf{Nonlinearity.} The system must exhibit nonlinear temporal evolution to evaluate the capability of PHYSGEN to learn complex dynamics. \item \textbf{Long-Term Sensitivity.} The system should demonstrate error accumulation during recursive forecasting. \item \textbf{Physical Interpretability.} The governing equations should provide measurable physical constraints. \item \textbf{Scalability.} The benchmark should allow extension from small systems to large graph-based representations. \end{itemize} %%%%%%%%%%%%%%%%%%%%%%%%%%%%%%%%%%%%%%%%%%%%%%%%%%%%%%%%%%%%%%%%%%%%%%%%%%%%%%%% \subsubsection{Lorenz-63 Chaotic System} %%%%%%%%%%%%%%%%%%%%%%%%%%%%%%%%%%%%%%%%%%%%%%%%%%%%%%%%%%%%%%%%%%%%%%%%%%%%%%%% The first benchmark is the Lorenz-63 system, a classical nonlinear chaotic dynamical system. It is defined by the following ordinary differential equations: \begin{equation} \frac{dx}{dt} = \sigma(y-x), \end{equation} \begin{equation} \frac{dy}{dt} = x(\rho-z)-y, \end{equation} \begin{equation} \frac{dz}{dt} = xy-\beta z . \end{equation} The system state is: \begin{equation} X(t) = [x(t),y(t),z(t)]^T . \end{equation} The Lorenz-63 benchmark evaluates: \begin{itemize} \item chaotic trajectory prediction; \item sensitivity to initial conditions; \item long-horizon stability; \item latent dynamics modeling. \end{itemize} The corresponding graph representation is: \begin{equation} G_{Lorenz} = (V,E,X), \end{equation} where each state variable is represented as a graph node. %%%%%%%%%%%%%%%%%%%%%%%%%%%%%%%%%%%%%%%%%%%%%%%%%%%%%%%%%%%%%%%%%%%%%%%%%%%%%%%% \subsubsection{Reaction-Diffusion System} %%%%%%%%%%%%%%%%%%%%%%%%%%%%%%%%%%%%%%%%%%%%%%%%%%%%%%%%%%%%%%%%%%%%%%%%%%%%%%%% The second benchmark represents spatially distributed nonlinear dynamics. The reaction-diffusion equation is defined as: \begin{equation} \frac{\partial u}{\partial t} = D\nabla^2u + R(u), \end{equation} where: \begin{itemize} \item $u(x,t)$ represents the physical field; \item $D$ is the diffusion coefficient; \item $R(u)$ represents the nonlinear reaction term. \end{itemize} A common formulation is the Gray-Scott system: \begin{equation} \frac{\partial u}{\partial t} = D_u\nabla^2u - uv^2 + F(1-u), \end{equation} \begin{equation} \frac{\partial v}{\partial t} = D_v\nabla^2v + uv^2 - (F+k)v. \end{equation} The spatial domain is discretized into interacting elements: \begin{equation} G_{RD} = (V_{grid},E_{neighbor},X). \end{equation} This benchmark evaluates: \begin{itemize} \item spatial message passing; \item local interaction modeling; \item conservation-aware evolution. \end{itemize} %%%%%%%%%%%%%%%%%%%%%%%%%%%%%%%%%%%%%%%%%%%%%%%%%%%%%%%%%%%%%%%%%%%%%%%%%%%%%%%% \subsubsection{Navier--Stokes Fluid Dynamics} %%%%%%%%%%%%%%%%%%%%%%%%%%%%%%%%%%%%%%%%%%%%%%%%%%%%%%%%%%%%%%%%%%%%%%%%%%%%%%%% The third benchmark considers incompressible fluid dynamics governed by the Navier--Stokes equations: \begin{equation} \frac{\partial u}{\partial t} + (u\nabla)u = -\frac{1}{\rho}\nabla p + \nu\nabla^2u, \end{equation} with the incompressibility condition: \begin{equation} \nabla\cdot u=0. \end{equation} The system state contains: \begin{equation} X = [u,v,p], \end{equation} where: \begin{itemize} \item $u,v$ are velocity components; \item $p$ is pressure. \end{itemize} The discretized fluid domain is represented as: \begin{equation} G_{NS} = (V_{cells},E_{flow},X). \end{equation} This benchmark evaluates: \begin{itemize} \item large-scale graph scalability; \item nonlinear spatio-temporal prediction; \item physical constraint preservation. \end{itemize} %%%%%%%%%%%%%%%%%%%%%%%%%%%%%%%%%%%%%%%%%%%%%%%%%%%%%%%%%%%%%%%%%%%%%%%%%%%%%%%% \subsubsection{Benchmark Summary} %%%%%%%%%%%%%%%%%%%%%%%%%%%%%%%%%%%%%%%%%%%%%%%%%%%%%%%%%%%%%%%%%%%%%%%%%%%%%%%% The benchmark suite is summarized in Table~\ref{tab:benchmark_systems}. \begin{table}[h] \centering \caption{Benchmark dynamical systems used for PHYSGEN evaluation.} \label{tab:benchmark_systems} \begin{tabular}{lccc} \hline \textbf{System} & \textbf{Type} & \textbf{Main Challenge} & \textbf{Evaluation Goal} \\ \hline Lorenz-63 & ODE & Chaotic evolution & Long-horizon stability \\ Reaction-Diffusion & PDE & Spatial nonlinear dynamics & Graph interaction learning \\ Navier--Stokes & PDE & Complex fluid dynamics & Scalability and physics \\ \hline \end{tabular} \end{table} %%%%%%%%%%%%%%%%%%%%%%%%%%%%%%%%%%%%%%%%%%%%%%%%%%%%%%%%%%%%%%%%%%%%%%%%%%%%%%%% \subsubsection{Benchmark Hierarchy} %%%%%%%%%%%%%%%%%%%%%%%%%%%%%%%%%%%%%%%%%%%%%%%%%%%%%%%%%%%%%%%%%%%%%%%%%%%%%%%% The selected systems form a progressive complexity hierarchy: \begin{equation} ODE \rightarrow PDE \rightarrow Large\text{-}scale\ Physical\ Simulation . \end{equation} This hierarchy enables systematic evaluation of PHYSGEN from fundamental nonlinear dynamics to complex scientific simulation scenarios. %%%%%%%%%%%%%%%%%%%%%%%%%%%%%%%%%%%%%%%%%%%%%%%%%%%%%%%%%%%%%%%%%%%%%%%%%%%%%%%% \subsubsection{Summary} %%%%%%%%%%%%%%%%%%%%%%%%%%%%%%%%%%%%%%%%%%%%%%%%%%%%%%%%%%%%%%%%%%%%%%%%%%%%%%%% The benchmark suite provides a comprehensive evaluation environment for PHYSGEN, covering chaotic systems, spatially distributed processes, and complex fluid dynamics. Together, these benchmarks allow assessment of prediction accuracy, physical consistency, stability, and scalability across different classes of nonlinear physical systems. \] 
